\documentclass{beamer}
\usepackage[utf8]{inputenc}

\usetheme{Madrid}
\usecolortheme{default}
\usepackage{amsmath,amssymb,amsfonts,amsthm}
\usepackage{txfonts}
\usepackage{tkz-euclide}
\usepackage{listings}
\usepackage{adjustbox}
\usepackage{array}
\usepackage{tabularx}
\usepackage{gvv}
\usepackage{lmodern}
\usepackage{circuitikz}
\usepackage{tikz}
\usepackage{graphicx}
\usepackage{mathtools}

\setbeamertemplate{page number in head/foot}[totalframenumber]

\usepackage{tcolorbox}
\tcbuselibrary{minted,breakable,xparse,skins}



\definecolor{bg}{gray}{0.95}
\DeclareTCBListing{mintedbox}{O{}m!O{}}{%
	breakable=true,
	listing engine=minted,
	listing only,
	minted language=#2,
	minted style=default,
	minted options={%
		linenos,
		gobble=0,
		breaklines=true,
		breakafter=,,
		fontsize=\small,
		numbersep=8pt,
		#1},
	boxsep=0pt,
	left skip=0pt,
	right skip=0pt,
	left=25pt,
	right=0pt,
	top=3pt,
	bottom=3pt,
	arc=5pt,
	leftrule=0pt,
	rightrule=0pt,
	bottomrule=2pt,
	toprule=2pt,
	colback=bg,
	colframe=orange!70,
	enhanced,
	overlay={%
		\begin{tcbclipinterior}
			\fill[orange!20!white] (frame.south west) rectangle ([xshift=20pt]frame.north west);
	\end{tcbclipinterior}},
	#3,
}
\lstset{
	language=C,
	basicstyle=\ttfamily\small,
	keywordstyle=\color{blue},
	stringstyle=\color{orange},
	commentstyle=\color{green!60!black},
	numbers=left,
	numberstyle=\tiny\color{gray},
	breaklines=true,
	showstringspaces=false,
}
\title{8.2.11}
\date{30th September, 2025}
\author{EE25BTECH11042-Nipun Dasari}

\begin{document}
	
	\frame{\titlepage}
	\begin{frame}{Question}
	Find the roots of the equation $x^2 + 3x - 10 = 0$ 
	\end{frame}
	
	\begin{frame}{Theoretical Solution}
	We express the equation as a parabola $y = x^2 + 3x - 10$, which can be written in the standard conic form $\vec{x}^\top\vec{V}\vec{x} + 2\vec{u}^\top\vec{x} + f = 0$ as:
	\begin{align}
		x^2 + 3x - y - 10 = 0
	\end{align}
	The components of the conic form are:
	\begin{align}
		\vec{V} = \myvec{1&0\\0&0}, \quad 2\vec{u} = \myvec{3\\-1} \implies \vec{u} = \myvec{3/2\\-1/2}, \quad f = -10
	\end{align}
	\end{frame}
	
	\begin{frame}{Theoretical Solution}
		We find the intersection with the x-axis, which has the parametric form $\vec{x} = \vec{h} + \kappa\vec{m}$, with a point on the line $\vec{h} = \myvec{0\\0}$ and direction vector $\vec{m} = \myvec{1\\0}$.
		
		The values of $\kappa$ for the intersection points  The solution is:
		\begin{align}
			\kappa = \frac{-\vec{m}^\top(\vec{V}\vec{h} + \vec{u}) \pm \sqrt{(\vec{m}^\top(\vec{V}\vec{h}+\vec{u}))^2 - (\vec{m}^\top\vec{V}\vec{m})(\vec{h}^\top\vec{V}\vec{h} + 2\vec{u}^\top\vec{h} + f)}}{\vec{m}^\top\vec{V}\vec{m}}
		\end{align}
		We will now calculate each term in this formula.
	\end{frame}
	
	\begin{frame}{Theoretical Solution}
		\textbf{Denominator Term:}
		\begin{align}
			\vec{m}^\top\vec{V}\vec{m} = \myvec{1 & 0}\myvec{1&0\\0&0}\myvec{1\\0} = \myvec{1 & 0}\myvec{1\\0} = 1
		\end{align}
		\textbf{Numerator Term 1:} Since $\vec{h} = \vec{0}$, this term is $-\vec{m}^\top\vec{u}$.
		\begin{align}
			-\vec{m}^\top\vec{u} = -\left(\myvec{1 & 0}\myvec{3/2\\-1/2}\right) = -\frac{3}{2}
		\end{align}
	\end{frame}
	
	\begin{frame}{Theoretical Solution}
		\textbf{Term under the Square Root:} Since $\vec{h} = \vec{0}$, this term is $(\vec{m}^\top\vec{u})^2 - (\vec{m}^\top\vec{V}\vec{m})(f)$.
		\begin{align}
			(\vec{m}^\top\vec{u})^2 &= \left(\frac{3}{2}\right)^2 = \frac{9}{4} \\
			(\vec{m}^\top\vec{V}\vec{m})(f) &= (1)(-10) = -10 \\
			\text{under root term} &= \frac{9}{4} - (-10) = \frac{9}{4} + \frac{40}{4} = \frac{49}{4}
		\end{align}
		Now, we substitute these calculated values back into the formula for $\kappa$:
	\end{frame}
	
	\begin{frame}{Theoretical Solution}
		\begin{align}
		\kappa = \frac{-3/2 \pm \sqrt{49/4}}{1} = -\frac{3}{2} \pm \frac{7}{2}
	\end{align}
	This gives two values for $\kappa$:
	\begin{align}
		\kappa_1 = -\frac{3}{2} + \frac{7}{2} = \frac{4}{2} = 2 \\
		\kappa_2 = -\frac{3}{2} - \frac{7}{2} = -\frac{10}{2} = -5
	\end{align}
	\end{frame}
	
	\begin{frame}{Theoretical Solution}
			The points of intersection are found using $\vec{x_i} = \vec{h} + \kappa_i\vec{m}$:
		\begin{align}
			\vec{x_1} = \myvec{0\\0} + 2\myvec{1\\0} = \myvec{2\\0} \\
			\vec{x_2} = \myvec{0\\0} - 5\myvec{1\\0} = \myvec{-5\\0}
		\end{align}
		The roots of the original equation are the x-coordinates of these intersection points.
		
		\textbf{The roots are 2 and -5.}
	\end{frame}
	

	\begin{frame}{Plot}
		\begin{figure}
			\centering
			\includegraphics[width=0.7\columnwidth]{../figs/Figure_1.png}
			\caption{}
			\label{fig:fig}
		\end{figure}
	\end{frame}
	
\end{document}